\bookchapter{บทที่ 56: ขัดเกลาและเผยแพร่ — ทดสอบ ดีบัก และขึ้น Google Play}

\chapterintro{ทำแอปเสร็จแล้ว ต่อไปล่ะ? บทนี้ตอบสามเรื่องของการจาก “รันได้” สู่ “ส่งมอบได้”: ดีบักหาบั๊กอย่างมีประสิทธิภาพยังไง เขียนเทสกันบั๊กย้อนกลับยังไง และเอาแอปขึ้น Google Play ให้คนทั้งโลกดาวน์โหลดยังไง การขึ้นสโตร์ไม่ใช่แค่พิธีกรรม — มันคือเส้นแบ่งระหว่างผลงานชิ้นหนึ่งกับสินค้าหนึ่งชิ้น}

\booksection{สามอาวุธดีบัก: Logcat, Breakpoint, ตัวตรวจ UI}

\begin{center}
\begin{tabularx}{\linewidth}{|X|X|X|}
\hline
\textbf{เครื่องมือ} & \textbf{ทำอะไร} & \textbf{ใช้ยังไง} \tabularnewline
\hline
Logcat & ดู log ของระบบ/ของเราที่พิมพ์เอง & ในโค้ดพิมพ์ Log.d("TAG", "ค่า=\$x") แล้วค้นหา TAG ในแผงเพื่อกรอง \tabularnewline
ดีบักด้วยเบรกพอยต์ & ให้โปรแกรมหยุดที่บรรทัดที่ต้องการ แล้วดูค่าตัวแปร & คลิกข้างเลขบรรทัดให้ขึ้นจุดแดง แล้วรันในโหมด Debug ไล่ทีละขั้น \tabularnewline
Layout Inspector & ดูโครงสร้างหน้าจอขณะแอปทำงานอยู่ & เมนู Tools Layout Inspector ดูแผนผังคอมโพเนนต์แบบสด ๆ \tabularnewline
\hline
\end{tabularx}
\end{center}

\begin{center}\small\color{gray}สามท่าเด็ดของการดีบัก\end{center}

\begin{tcbinfo}{เคล็ดลับ: วิธีอ่านข้อความ error?}
เมื่อแอปพัง ใน Logcat จะมี stack trace สีแดง (Stacktrace) ให้ไล่จากล่างขึ้นบน หาบรรทัดแรกที่ขึ้นต้นด้วยชื่อแพ็กเกจของเราเอง (เช่น com.example.yourapp) บรรทัดนั้นมักเป็น “ที่เกิดเหตุ” พอดี ดู “ชนิดของ Exception” ก่อน (null pointer? แปลงชนิดข้อมูลผิด?) แล้วค่อยดู “at ... ไฟล์ของเรา:บรรทัด” เทคนิคนี้แก้ crash ได้ถึง 80\%

\end{tcbinfo}

\booksection{การทดสอบ: ทำประกันให้โค้ด}

การทดสอบด้วยมือไม่น่าเชื่อถือ — แก้ฟีเจอร์นี้ไป อีกฟีเจอร์ก็พังเงียบ ๆ ไปด้วย การทดสอบคือทั้ง “ยาถอนพิษ” และ “เสื้อเกราะกันกระสุน” การทดสอบใน Android แบ่งเป็น 3 ชั้น มือใหม่เริ่มจากทำความเข้าใจชั้นแรกให้แน่นก่อน:

\begin{center}
\begin{tabularx}{\linewidth}{|X|X|X|X|}
\hline
\textbf{ชั้น} & \textbf{ทดสอบอะไร} & \textbf{วางไว้ที่ไหน} & \textbf{ความเร็ว} \tabularnewline
\hline
ยูนิตเทสต์ & ลอจิกล้วน ๆ (การคำนวณ การตรวจสอบ) & app/src/test/ & ระดับวินาที \tabularnewline
อินสตรูเมนต์เทสต์ & พฤติกรรมหน้าจอบนอุปกรณ์จริง & app/src/androidTest/ & ระดับนาที \tabularnewline
เทสต์ด้วยมือ & ให้คนจริงสัมผัสและใช้งานจริง & —— & ต้องมีตลอดไป \tabularnewline
\hline
\end{tabularx}
\end{center}

\begin{center}\small\color{gray}การทดสอบสามชั้น\end{center}

\begin{lstlisting}[language=Kotlin]
// แยกตรรกะ "แปลงข้อความจำนวนเงินเป็นตัวเลข" ออกเป็นฟังก์ชันล้วน
fun parseAmount(input: String): Double? {
    val v = input.trim().toDoubleOrNull()
    return if (v != null && v > 0) v else null
}

// ทดสอบมัน (JUnit4)
@Test
fun emptyInputReturnsNull() {
    assertEquals(null, parseAmount(""))
}
@Test
fun negativeReturnsNull() {
    assertEquals(null, parseAmount("-5"))
}
@Test
fun normalValueParses() {
    assertEquals(12.5, parseAmount("12.5"))
}
\end{lstlisting}

การตรวจสอบจำนวนเงินและการคำนวณสถิติในแอปบันทึกรายรับ-รายจ่าย ควรมียูนิตเทสต์คอยปกป้องไว้ วิธีรันเทสต์: คลิกขวาที่โฟลเดอร์ test Run ถ้าขึ้นเขียวผ่านหมด ก็แก้โค้ดได้อย่างมั่นใจ

\booksection{การเซ็นและแพ็กเกจ: ประทับตราให้แอป}

Android กำหนดว่าไฟล์ติดตั้งทุกไฟล์ต้องถูกเซ็นก่อน — เหมือนประทับตราทางการ ยืนยันว่า “แพ็กเกจนี้มาจากเราและไม่ถูกแก้ไข” ตอนดีบัก Android Studio จะเซ็นด้วย debug key ให้อัตโนมัติ ส่วนการเผยแพร่ต้องใช้ release key ของเราเอง ขั้นตอนมีดังนี้:

\begin{enumerate}
  \item เมนู Build Generate Signed Bundle / APK;
  \item สร้างไฟล์ keystore (คลังกุญแจ) — ตั้งรหัสผ่านและเก็บรักษาให้ดี;
  \item เลือก release + รูปแบบ AAB แล้วบิลด์เป็นอันเสร็จ.
\end{enumerate}

\textbf{ระวังพลาด: keystore หาย = แอปหาย}\par

หลังขึ้นสโตร์ การอัปเดตแอปต้องใช้ลายเซ็นชุดเดิมเสมอ ไม่เช่นนั้นสโตร์จะปฏิเสธการอัปเดต ถ้าไฟล์ keystore และรหัสผ่านหายไป คุณจะเหลือทางเดียวคือสร้างแอปใหม่ขึ้นมาขึ้นสโตร์ใหม่ตั้งแต่ต้น — ผู้ใช้เดิมทั้งหมดเชื่อมต่อไม่ได้เลย อย่าลืมสำรอง keystore และรหัสผ่านไว้ในที่ปลอดภัย (เข้ารหัสในคลาวด์ + สมุดจด) นี่คือบทเรียนที่แพงที่สุดบทหนึ่งของนักพัฒนา Android

\booksection{กระบวนการครบวงจรของการขึ้น Google Play}

สมัครบัญชีนักพัฒนา

เตรียมสื่อสำหรับสโตร์

อัปโหลด AAB

กรอกตารางความปลอดภัยข้อมูล

รีวิว (หลายชั่วโมง\textasciitilde{}หลายวัน)

เผยแพร่บนสโตร์

\begin{center}\small\color{gray}ภาพที่ 6: หกขั้นตอนสู่ Google Play\end{center}

\begin{itemize}
  \item ลงทะเบียน: ค่าสมัครครั้งเดียว 25 ดอลลาร์ ต้องมีบัญชี Google ที่ใช้งานได้;
  \item สื่อ: ไอคอนแอป (512×512), ภาพหน้าจอ (อย่างน้อย 2 ภาพ), คำอธิบายฟีเจอร์, ลิงก์นโยบายความเป็นส่วนตัว (ถ้ามีการเก็บข้อมูลต้องใส่);
  \item แบบฟอร์มความปลอดภัยของข้อมูล: แจ้งตรง ๆ ว่าเก็บข้อมูลอะไรบ้าง — แอปบันทึกรายรับ-รายจ่ายที่ไม่ต่อเน็ตเลือก “ไม่เก็บเลย” ก็ได้;
  \item ตรวจสอบ: บัญชีใหม่ที่ขึ้นสโตร์ครั้งแรกมักถูกตรวจโดยคน ใช้เวลาตั้งแต่ไม่กี่ชั่วโมงถึงหลายวัน ถ้าถูกส่งกลับให้แก้ตามข้อเสนอแนะ.
\end{itemize}

สำหรับผู้ใช้ในจีน: Google Play ต้องใช้สภาพแวดล้อมเครือข่ายพิเศษ และการเผยแพร่ในจีนแผ่นดินใหญ่มักทำผ่านสโตร์ในประเทศ (Huawei, Xiaomi, Yingyongbao ฯลฯ) ขั้นตอนขึ้นสโตร์ในประเทศคล้ายกัน แต่แต่ละสโตร์มีเกณฑ์ตรวจสอบแบบของตัวเอง — AAB/APK ชุดเดียวกันเผยแพร่หลายสโตร์ได้ สำหรับการฝึกมือ ขึ้นได้สักสโตร์หนึ่งก็ถือว่าชนะแล้ว

\booksection{ทำสื่อในสโตร์อย่างไร: ความประทับใจแรกก็คือยอดดาวน์โหลด}

หน้าสโตร์คือหน้าต่างโชว์ที่ผู้ใช้เห็นตอน “เดินเลือกซื้อ” สื่อ 3 อย่างตัดสินความประทับใจแรก:

\begin{itemize}
  \item ไอคอน: 512×512 อย่าใช้ภาพหน้าจอมาแก้ขัด ถ้าไม่มีทักษะออกแบบ ให้ใช้เครื่องมือ AI สร้างรูปทรงเรขาคณิตเรียบ ๆ ขึ้นมา เกณฑ์ผ่านคือหดเล็กลงบนหน้าจอหลักแล้วยังมองปุ๊บรู้เรื่องทันที;
  \item ภาพหน้าจอ: อย่างน้อย 2 ภาพ ดีที่สุด 4-6 ภาพ โชว์สถานการณ์ใช้งานที่สำคัญที่สุด แต่ละภาพแปะประโยคสั้น ๆ เป็นภาษาคน (“จดรายการเสร็จในวินาทีเดียว” ดีกว่า “โซลูชันบันทึกบัญชีที่มีประสิทธิภาพ” สิบเท่า);
  \item คำอธิบาย: ประโยคแรกต้องบอกให้ชัดว่า “นี่คืออะไร + ทำไว้ให้ใครใช้” อย่าเรียงร้อยคำว่า “เป็นมืออาชีพที่สุด เก่งที่สุด” — ผลการค้นหาในสโตร์ขึ้นกับคีย์เวิร์ดและอัตราการดาวน์โหลดเป็นหลัก ความจริงใจได้ผลกว่าการอวด.
\end{itemize}

\begin{tcbinfo}{เคล็ดลับ: เปิดเทียบแอปเพื่อนร่วมอาชีพ}
ก่อนขึ้นสโตร์ หมั่นเข้าไปดูแอปประเภทเดียวกันใน 10 อันดับแรก จับภาพมาเปรียบเทียบ: ไอคอนเขาสไตล์ไหน? ภาพหน้าจอขายจุดไหน? เขียนชื่ออย่างไร? เลียนแบบไม่ใช่ลอก — เรียนรู้สูตร แล้วใส่เนื้อหาของตัวเอง นี่คือบทเรียนที่ทำให้ผลิตภัณฑ์ใหม่เริ่มต้นได้เร็วที่สุด

\end{tcbinfo}

\booksection{เช็กลิสต์เล็ก ๆ เรื่องประสิทธิภาพกับการรายงาน crash}

ก่อนขึ้นสโตร์ สองเรื่องนี้คุ้มค่าที่จะลงเวลาอย่างละครึ่งชั่วโมง:

\booksubsection{เช็กประสิทธิภาพด้วยตัวเอง 4 ข้อ}

\begin{enumerate}
  \item ลิสต์ใช้ LazyColumn: รายการยาว ๆ ห้ามใช้ Column เรียงตรง ๆ เด็ดขาด — lazy loading คือบทเรียนแรกเรื่องประสิทธิภาพบน Android;
  \item อย่ายัดภาพต้นฉบับลงไปตรง ๆ: ใช้ไลบรารีภาพ Coil ย่อขนาดและแคชให้อัตโนมัติ ภาพต้นฉบับ 5MB ไฟล์เดียวก็พอทำให้เครื่องราคาถูกกระตุกจนตาย;
  \item หลีกเลี่ยงการทำงานหนักบน main thread: ฐานข้อมูล เครือข่าย ไฟล์ใหญ่ ทำใน coroutine ทั้งหมด;
  \item ใช้ Profiler วัดดู: เครื่องมือวิเคราะห์ประสิทธิภาพในตัวของ Android Studio บอกได้ทั้งหน่วยความจำ CPU และการกินแบต — วัดก่อน แล้วค่อยปรับปรุง อย่าเดาด้วยความรู้สึก.
\end{enumerate}

\booksubsection{รายงาน crash}

แอปไม่พังตอนเราทดสอบ ไม่ได้แปลว่าจะไม่พังบนเครื่องผู้ใช้ ก่อนขึ้นสโตร์ ต่อ Firebase Crashlytics (ของ Google เอง ฟรี): แค่ไม่กี่บรรทัดของการตั้งค่า เมื่อแอปของผู้ใช้พัง มันจะส่ง stack trace ขึ้นมาให้อัตโนมัติ คุณเห็นในหลังบ้านเลยว่า “รุ่นไหน เวอร์ชันไหน พังบรรทัดไหน” การขึ้นสโตร์โดยไม่มีระบบเฝ้าดู crash เท่ากับขับรถล้อตา

\begin{tcbpractice}{ลงมือทำ: “ส่งมอบ” แอปบันทึกรายรับ-รายจ่ายสักครั้ง}
(1) เติมยูนิตเทสต์ 3 ตัวให้ลอจิกอย่าง parseAmount และอื่น ๆ; (2) ใช้ Logcat จับ crash ที่เราสร้างขึ้นเองด้วยมือตัวเองแล้วแก้ให้หาย; (3) สร้าง APK ที่เซ็นแบบ release ส่งให้เพื่อนติดตั้ง (นี่คือ “พิธีกรรม” ของผลงาน); (4) ถ้ามีแรงเหลือ ลงทะเบียนบัญชีนักพัฒนาแล้วยื่นแอปบันทึกรายรับ-รายจ่ายขึ้นสโตร์จริง ๆ สามข้อแรกทำเสร็จได้ในคืนนี้

\end{tcbpractice}

\begin{tcbsummary}{สรุปบทนี้}
\begin{itemize}
  \item สามท่าเด็ดของการดีบัก: Logcat ดู log, เบรกพอยต์ดูตัวแปร, Layout Inspector ดูหน้าจอ
  \item อ่าน crash: ไล่ stack จากล่างขึ้นบน หาบรรทัดแรกที่เป็นชื่อแพ็กเกจของเรา
  \item ทดสอบ 3 ชั้น: ยูนิต (ระดับวินาที เริ่มจากอันนี้) / อินสตรูเมนต์ / แมนนวล
  \item เผยแพร่ต้องเซ็น: สำรอง keystore กับรหัสผ่านให้ดี ถ้าหาย แอปก็ต้อง “เปลี่ยนตัว” ไป
  \item หกขั้นตอนขึ้นสโตร์: ลงทะเบียน เตรียมสื่อ อัปโหลด AAB กรอกแบบฟอร์มความปลอดภัยของข้อมูล ตรวจสอบ เผยแพร่
\end{itemize}

\end{tcbsummary}
